\documentclass[11pt,a4paper]{article}
\usepackage[margin=2.2cm]{geometry}
\usepackage[T1]{fontenc}
\usepackage{amsmath}
\usepackage{booktabs}
\usepackage{tabularx}
\usepackage{array}
\usepackage{hyperref}

\newcolumntype{L}{>{\raggedright\arraybackslash}X}

\title{Report of Changes:\\ Two Extensions of Fly-CL}
\author{}
\date{13 October 2026}

\begin{document}
\maketitle

\noindent Base code: official Fly-CL repository (\url{https://github.com/gfyddha/Fly-CL}),
upstream commit \texttt{193b1b8}.
Project tasks: (1) Activity-Dependent Learnable Sparse Projection,
(2) Variance-Adaptive Generalized Ridge.

\section{Principle followed}

\begin{itemize}
  \item \textbf{Minimal changes to the original code.} Every new behaviour is behind a
        command-line flag whose default reproduces the original Fly-CL. New logic is in new files.
  \item \textbf{Every step is a git commit}, so the history shows what changed and when.
\end{itemize}

Table~\ref{tab:size} gives the size of the change to the original files
(\texttt{git diff 193b1b8 HEAD -{}-numstat}).

\begin{table}[h]
\centering
\caption{Size of the change per file.}
\label{tab:size}
\begin{tabular}{lrrl}
\toprule
File & Lines added & Lines removed & Kind \\
\midrule
\texttt{main.py}       & 75 & 5 & modified \\
\texttt{utils.py}      & 16 & 2 & modified \\
\texttt{extensions.py} & 60 & 0 & new (both methods) \\
\texttt{.gitignore}    & 6  & 1 & modified \\
\texttt{readme.md}     & 9  & 0 & modified \\
\bottomrule
\end{tabular}
\end{table}

In \texttt{main.py} only four original statements were touched: the \texttt{utils} import, the two
\texttt{feature\_extract(...)} calls (one extra argument) and the Cholesky line (kept verbatim,
moved into an \texttt{else} branch). The folders \texttt{datasets/}, \texttt{models/} and
\texttt{scripts/} are untouched.

\section{Git history}

\begin{table}[h]
\centering
\caption{Commits on top of the upstream code (all on 2026-10-13).}
\begin{tabularx}{\textwidth}{llL}
\toprule
Commit & Time & What \\
\midrule
\texttt{6ee6341} & 17:13 & Experiment infrastructure: feature cache, BWT / last-stage metrics, json output \\
\texttt{12e121c} & 18:04 & Task 1: activity-dependent learnable sparse projection (Hebbian) \\
\texttt{d5bf460} & 18:04 & Task 2: variance-adaptive generalized ridge \\
\texttt{fb04d49} & 19:00 & Experiment, validation-sweep and analysis scripts \\
\texttt{d8a03e9} & 19:01 & Results, analysis figures and report \\
\bottomrule
\end{tabularx}
\end{table}

Useful commands: \texttt{git log -{}-stat}, \texttt{git show <commit>},
\texttt{git diff 193b1b8 HEAD -{}- main.py}.

\section{Change 1: Experiment infrastructure (commit \texttt{6ee6341})}

No change to the learning algorithm.

\begin{table}[h]
\centering
\begin{tabularx}{\textwidth}{>{\raggedright\arraybackslash}p{3.3cm}LL}
\toprule
Where & Change & Why \\
\midrule
\texttt{utils.py}: \texttt{feature\_extract} &
optional \texttt{cache\_path} argument: load features from disk if present, otherwise extract and save &
The backbone is frozen, so features are identical for all methods. Caching makes all methods run
on exactly the same inputs and cuts a run from about 7 min to about 1.5 min. \\
\texttt{utils.py}: \texttt{save\_results} & new 4-line helper writing a json file &
Results are aggregated by script instead of copied from logs. \\
\texttt{main.py} flags & \texttt{-{}-feature\_cache}, \texttt{-{}-save\_json}, \texttt{-{}-save\_state} &
Switch the above on; all default to \texttt{None} (off). \\
\texttt{main.py} end of file & prints last-stage accuracy, BWT, training time excluding feature extraction &
The project asks for these metrics; the original printed only per-stage and accumulated accuracy and total time. \\
\texttt{main.py} timer & \texttt{torch.cuda.synchronize()} before the training timer stops &
GPU calls are asynchronous; without it the measured time can miss pending work. \\
\texttt{main.py} bookkeeping & \texttt{ridges} list, \texttt{initial\_projection\_matrix}, \texttt{ridge\_diag} &
Needed only for the analysis (GCV choice, mask check, conditioning). \\
\bottomrule
\end{tabularx}
\end{table}

Metric definitions added. Let $R_{t,i}$ be the accuracy on task $i$ after learning task $t$,
with $T$ tasks:
\begin{align*}
\text{Last-stage accuracy} &= \text{accuracy on all classes after task } T, \\
\text{BWT} &= \frac{1}{T-1}\sum_{i=1}^{T-1}\left(R_{T,i} - R_{i,i}\right).
\end{align*}

\section{Change 2: Task 1, Activity-Dependent Learnable Sparse Projection (commit \texttt{12e121c})}

New function \texttt{hebbian\_update(projection\_matrix, features, lr, epochs, hebb\_k, delta)}
in \texttt{extensions.py}.

\begin{table}[h]
\centering
\caption{How each instruction of Task 1 is met.}
\begin{tabularx}{\textwidth}{LL}
\toprule
Instruction & Implementation \\
\midrule
Fix a binary connectivity mask at initialization &
$M = [W \neq 0]$ taken from the original random projection; all updates are multiplied by $M$. \\
Replace fixed values by weights learned with a local Hebbian / anti-Hebbian rule &
BioHash-style rule on the active connections, Eq.~\eqref{eq:hebb}. \\
Keep the projection sparsity and the same Top-$k$ &
The returned matrix has the same non-zero pattern; the Top-$k$ code in \texttt{main.py} is unchanged. \\
Stable variant: adapt on the first task, then freeze & \texttt{-{}-projection hebbian\_frozen} \\
Optional: continual updates & \texttt{-{}-projection hebbian\_continual} \\
\bottomrule
\end{tabularx}
\end{table}

For unit $\mu$ with weight row $W_\mu$, mask row $M_\mu$ and input $x$:
\begin{equation}
\Delta W_\mu = g_\mu \left( M_\mu \odot x - \langle W_\mu, x\rangle\, W_\mu \right),
\qquad
g_\mu =
\begin{cases}
1 & \text{if } \mu \text{ is among the } k_h \text{ most active units for } x,\\
-\delta & \text{if } \mu \text{ is among the next } k_h \text{ units},\\
0 & \text{otherwise.}
\end{cases}
\label{eq:hebb}
\end{equation}
Features are centred and $\ell_2$-normalised for learning only. The update of each unit is divided
by the number of samples that drove it (activity-dependent step). Each row is re-normalised to
the norm it had at initialization, so the expanded features keep the original scale and the ridge
grid needs no change.

Change in \texttt{main.py} (3 lines in the task loop, plus 5 flags):
\begin{verbatim}
if args.projection == 'hebbian_continual' or \
        (args.projection == 'hebbian_frozen' and task == 0):
    projection_matrix = hebbian_update(projection_matrix, train_embeddings,
                                       args.hebb_lr, args.hebb_epochs,
                                       args.hebb_k, args.hebb_delta)
\end{verbatim}
(The line breaks above are for the page; in the file the condition is on one line.)

New flags: \texttt{-{}-projection \{fixed,hebbian\_frozen,hebbian\_continual\}} (default \texttt{fixed}),
\texttt{-{}-hebb\_lr} (0.1), \texttt{-{}-hebb\_epochs} (1), \texttt{-{}-hebb\_k} (0.01),
\texttt{-{}-hebb\_delta} (0.0).
The Hebbian step is inside the training timer, so its cost is included in the reported time.

\section{Change 3: Task 2, Variance-Adaptive Generalized Ridge (commit \texttt{d5bf460})}

New function \texttt{adaptive\_ridge\_diag(G, ridge, alpha, clamp=100)} in \texttt{extensions.py}.

\begin{table}[h]
\centering
\caption{How each instruction of Task 2 is met.}
\begin{tabularx}{\textwidth}{LL}
\toprule
Instruction & Implementation \\
\midrule
Use the diagonal of the accumulated Gram matrix as activity / variance &
\texttt{activity = torch.diagonal(G)} \\
Replace $\lambda I$ by $\Lambda = \mathrm{diag}(\lambda_1,\dots,\lambda_m)$ &
returns the vector $\lambda_j$; \texttt{main.py} adds \texttt{torch.diag(ridge\_diag)} to $G$ \\
Dimension-specific coefficients from streaming statistics, GCV coefficient as global scale &
Eq.~\eqref{eq:ridge} \\
Keep streaming updates of $G$, $Q$; efficient positive-definite solver &
$G$, $Q$, GCV untouched; solved with \texttt{torch.linalg.cholesky} and \texttt{cholesky\_solve} \\
\bottomrule
\end{tabularx}
\end{table}

\begin{equation}
\lambda_j = \lambda_{\mathrm{GCV}}\; s_j,
\qquad
s_j = \left(\frac{G_{jj}}{\tfrac{1}{m}\sum_i G_{ii}}\right)^{\alpha},
\label{eq:ridge}
\end{equation}
with $s_j$ clamped to $[0.01, 100]$ and normalised to mean 1. The classifier is
$W_o = (G + \Lambda)^{-1} Q$.
$\alpha = 0$ gives the original scalar ridge; $\alpha = 1$ equals ridge on variance-standardised
dimensions; $\alpha > 0$ regularizes highly active dimensions more.

Change in \texttt{main.py} (the original line is kept in the \texttt{else} branch):
\begin{verbatim}
if args.ridge_type == 'adaptive':
    ridge_diag = adaptive_ridge_diag(G, ridge.item(), args.ridge_alpha)
    L = torch.linalg.cholesky(G + torch.diag(ridge_diag))
else:
    ridge_diag = ridge * torch.ones(G.size(dim=0)).to(device)
    L = torch.linalg.cholesky(G + ridge * torch.eye(G.size(dim=0)).to(device))
\end{verbatim}

New flags: \texttt{-{}-ridge\_type \{scalar,adaptive\}} (default \texttt{scalar}),
\texttt{-{}-ridge\_alpha} (1.0).

\section{Change 4: Experiment and analysis scripts (commit \texttt{fb04d49})}

All new, in \texttt{experiments/}; none of them changes the method.

\begin{table}[h]
\centering
\begin{tabularx}{\textwidth}{lL}
\toprule
File & Purpose \\
\midrule
\texttt{run.sh} & One CIFAR-100 run with the configuration of \texttt{scripts/test\_cifar.sh}, cached features, json + log. \\
\texttt{make\_val\_cache.py} & Holds out the last 10\% of each task's training features as a validation split, so hyperparameters are not chosen on the test set. \\
\texttt{sweep\_val.sh} & The validation sweep for the Hebbian and ridge hyperparameters. \\
\texttt{run\_all.sh} & The final comparison: 8 configurations $\times$ 3 seeds. \\
\texttt{analyze.py} & Prototype correlation matrices before / after projection, mask check, condition numbers, per-dimension ridge analysis, figures. \\
\texttt{summarize.py} & Aggregates all json files into \texttt{results/summary.md}. \\
\bottomrule
\end{tabularx}
\end{table}

\section{Change 5: Results and documentation (commit \texttt{d8a03e9})}

\begin{itemize}
  \item \texttt{results/} (test-set runs, summary, figures), \texttt{results\_val/} (validation sweep).
  \item \texttt{RESULTS.md} (report), \texttt{CHANGELOG.md} (running log), a short section in \texttt{readme.md}.
  \item \texttt{.gitignore}: \texttt{.venv/}, \texttt{cache/}, \texttt{data/}, \texttt{*.pt} (large or machine-specific files).
\end{itemize}

\section{How to verify that the original behaviour is preserved}

\begin{itemize}
  \item With default flags no new branch is executed: \texttt{-{}-projection fixed},
        \texttt{-{}-ridge\_type scalar}, no cache, no json.
  \item The baseline run gave 93.21 average accuracy before the extension code was added and 93.20
        after (seed 1993). A difference of 0.01 corresponds to about one test image at one stage;
        its cause was not investigated (most likely run-to-run GPU floating-point variation).
  \item \texttt{experiments/analyze.py} asserts that the learned projection has exactly the non-zero
        pattern of the initial one and that all methods start from the same initial matrix.
\end{itemize}

\section{Environment notes (not code changes)}

\begin{itemize}
  \item Runs used torch 2.14.1, torchvision 0.29.1, numpy 2.5 instead of the versions in the readme;
        the original code runs on them unmodified. \texttt{timm==0.9.16} needs
        \texttt{huggingface\_hub==0.23.5}.
  \item CIFAR-100 was copied into \texttt{data/} with the official md5 checksums verified.
  \item One row of the original random projection has 298 instead of 300 non-zeros for these seeds,
        as generated by the unmodified initialization; it is the same for all methods.
\end{itemize}

\section{How to run}

\begin{verbatim}
# original Fly-CL
python main.py --dataset CIFAR-100 ...
# Task 1 (stable variant)
python main.py --dataset CIFAR-100 ... --projection hebbian_frozen
# Task 2
python main.py --dataset CIFAR-100 ... --ridge_type adaptive --ridge_alpha 1.0
# everything used in the report
experiments/run_all.sh
\end{verbatim}

\end{document}
